This study shows that a sequence model's incremental performance is not determined by the model
alone. When the same source peaks were used and the model class, fold assignment and estimator were held fixed, changing
the construction of the negative set altered the estimated contribution of a 4-mer model by
5.42-fold under the conventional two-stage estimator and by 4.84-fold under the primary
cross-fitted estimator. The corresponding two-stage spans ranged from 3.7- to 7.4-fold across
three model classes. Within the study panel, the protocol yielding the lowest apparent AUROC produced
the largest increment over composition, whereas the protocol yielding the highest apparent
AUROC produced the smallest increment. Negative-set construction therefore changes both the
difficulty of the task and the quantity being measured; the two should not be treated as
interchangeable indicators of stringency.

Five reporting practices follow from these results.

\begin{enumerate}\setlength{\itemsep}{2pt}
\item \textbf{Describe how negatives were constructed.} The candidate pool, matching variables,
  tolerances, achieved match quality and unconstrained covariates should be reported in enough
  detail to reproduce the set. Nominal tolerances alone did not capture the achieved matching
  in the study data.
\item \textbf{Report the composition-only AUROC under the same protocol.} This baseline is
  inexpensive to fit and provides a direct measure of the discrimination available from
  composition. In the bias-aware arm, the 19-feature composition model outperformed the 4-mer
  in 53 of 94 datasets.
\item \textbf{State the order of the composition baseline.} Incremental performance is defined
  relative to the features already represented. Raising the baseline from dinucleotide to
  trinucleotide composition reduced every contribution, although protocol dependence remained.
\item \textbf{Confirm that the baseline remains estimable as its order increases.} In this panel,
  out-of-fold baseline performance improved through order three but not at order four. Beyond that
  point, apparent increments can reflect estimation error in the baseline rather than additional
  information in the sequence model.
\item \textbf{Calibrate the estimator with a null model.} A 2-mer model provides such a control
  because its information is already contained in the dinucleotide baseline. The conventional
  estimator nevertheless returned $+0.0111$ to $+0.0137$ (Section~\ref{sec:floor}), showing that
  increments of this magnitude cannot be interpreted without an estimator-specific null.
\end{enumerate}

These recommendations imply that contributions estimated under different negative-set
protocols should not be compared directly. Each protocol defines a different classification
problem, and an increment is meaningful only relative to the baseline and observations used to
estimate it.

The bias-aware protocol illustrates why this distinction matters. Drawing negatives from sites
bound by other RBPs can reduce broad accessibility and composition artefacts, but it also poses
a biologically different task from separating bound sites from genomic background. RBPs differ
in their preferred transcript regions, and transcript regions have characteristic nucleotide
composition. Stratifying the donor draw by region reduced the bias-aware composition baseline
by 0.020 AUROC, accounting for 46\% of its excess over the GC arm, and decreased its measured
contribution by approximately one quarter. The arm nevertheless remained the lowest of the
three in incremental performance. Thus, transcript region explains part, but not all, of the
bias-aware baseline. This is not an argument against bias-aware negatives; it shows that their
results answer a different question and require protocol-specific interpretation.

No standard reparameterisation made the contributions comparable across protocols. Dividing by
the baseline's remaining headroom reduced the span, but the fitted exponent depended on the
aggregation rule, benchmark and model class. It therefore cannot serve as a transferable
normalisation. The protocol ordering was also preserved for likelihood-ratio statistics,
McFadden's $R^2$, average precision and integrated discrimination improvement. The magnitude
was scale-dependent: for the 4-mer, the three-protocol span was 5.42-fold in AUROC and
approximately two-fold on the unbounded scales. The qualitative dependence is consequently not
an artefact of the AUROC ceiling, although the headline fold change is specific to the AUROC
scale.

The baseline order introduces a second source of dependence. A 4-mer model evaluated against a
trinucleotide baseline largely measures the remaining fourth-order composition rather than a
general notion of motif recognition. Increasing the baseline order is therefore informative but
not neutral: it is structurally less favourable to $k$-mer models than to models that can use
position or longer-range context. I interpret these analyses as a characterisation of what the
baseline absorbs, not as a model-capacity ranking.

These findings extend the observation by \citet{horlacher2023} that negative-set construction
changes apparent RBP-prediction performance. The present analysis shows that the same choice also
changes incremental performance relative to an explicit baseline. In the study panel, the magnitude
of protocol-associated variation was comparable to that produced by changing model class:
within-protocol model spans were 2.62- to 4.14-fold, whereas within-model protocol spans were
3.7- to 7.4-fold. Similar sensitivity has been reported in transcription-factor binding
prediction \citep{tourne2026}, suggesting that the issue may extend beyond eCLIP, although the
current study does not test that generalisation directly.

\subsection*{Limitations}

The estimated magnitudes are conditional on the 19-feature mono- and dinucleotide baseline.
Every contribution decreased when the baseline was extended to trinucleotides, while the
protocol ordering persisted. The existence of protocol dependence was therefore robust to the
tested baseline orders, but its size was not. Neither the present data nor a transformation can
define a protocol-independent contribution.

The conventional two-stage estimator introduced an additional limitation. For a 2-mer model
whose information is contained in the baseline, it produced a positive contribution rather than
the expected zero. The source was an information route between the inner score generation and
outer evaluation folds. Cross-fitting the score covariate reduced the floor by at least 95\%
in every arm and narrowed the 4-mer span from 5.42-fold to 4.84-fold without removing it. This
correction was not run for the convolutional network or SpliceBERT because it would require new
neural training. Their spans are therefore exploratory two-stage estimates and should not be
compared numerically with the primary cross-fitted 4-mer result.

External validation supports the general dependence on negative-set construction but narrows
its interpretation. In 135 datasets from \citet{horlacher2023} that were absent from the study panel,
the two negative sets produced a 1.69-fold two-stage span (95\% CI 1.41--2.03) and a 1.66-fold
cross-fitted span (1.40--1.98). Chromosome-blocked refolding gave corresponding estimates of
1.73 (1.43--2.10) and 1.67 (1.40--2.00). These datasets did not overlap the study datasets, but 31 of the
108 proteins overlapped;
excluding all shared proteins still gave a 1.66-fold span. The analysis is therefore external in
dataset construction and processing, not in assay, organism or all biological entities. It
remains ENCODE eCLIP from K562 and HepG2 cells.

The inverse relation between apparent difficulty and incremental contribution did not replicate
externally. In the Horlacher benchmark, the more difficult arm also yielded the smaller
contribution, and the per-dataset association was null. The transferable result is therefore
that negative-set construction changes the measured contribution, not that apparent AUROC and
contribution must move in opposite directions. The external benchmark was also known before the
held-out subset was scored. Its criteria were specified before analysis of that subset, but the
study should not be interpreted as a prospective search across independent benchmarks.

The study covers three model classes, two cell lines and one assay. The 4-mer model, small
convolutional network and 19.78~M-parameter SpliceBERT form a capacity ladder rather than a
survey of published RBP predictors. No model larger than 20~M parameters was tested. The
cross-protocol training and evaluation analysis (Section~\ref{sec:transport}) was limited to the
4-mer; extending it to the neural models would require retraining. Its off-diagonal cells also combine a protocol change
with distribution shift, so the estimated training and evaluation shares are descriptive and
not causal.

Several design choices were fixed before analysis. All primary arms used a 1:1 class ratio,
whereas genome-wide applications are much more imbalanced. AUROC is prevalence-invariant, but
precision-based measures are not. Sensitivity analyses at 1:2, 1:4 and 2:1 preserved the
protocol ordering, although the magnitudes varied. Window length was fixed at 101~nt and was not
systematically varied. A longer window would provide more sequence for estimating composition
and could alter both baseline and incremental performance.

Chromosome-blocked folds prevent loci from being divided across training and test sets, but they
do not eliminate similarity between homologous loci on different chromosomes. Only 176 of
2.7 million windows belonged to gene names represented in more than one fold, primarily
pseudo-autosomal and multicopy small-RNA genes. Gene-clustered refitting changed the two-arm
contrast by $0.0003$. Three alternative chromosome assignments varied it by $0.00058$. These
results indicate limited sensitivity to the fold assignment for the 4-mer, but the neural models
were evaluated only on the frozen chromosome partition.

Window centring had a larger effect. The discriminative signal occurred a median of 23~nt from
the peak midpoint, and summit-centred windows could not be tested because the required summit
field was absent from the examined peak files. Across midpoint, transcript-oriented boundary
and displaced-midpoint definitions, the two-arm contrast varied by $0.0077$ in ten datasets.
The published midpoint definition produced the largest of the three estimates, so window
centring remains an important source of uncertainty.

Negative construction is stochastic. Independent redraws changed panel means less than the
between-protein sampling uncertainty: incorporating draw variability increased interval widths
by 0.49--0.81\% in the composition-matched arms and by 3.6\% in the bias-aware arm. The
ordering of protocols was unchanged. For the composition-matched arms, redraws condition on the
positive windows retained by the original matching pass and therefore quantify variability in
negative selection rather than the complete construction process.

Neural-model initialisation was not seeded across the 2830 committed fold runs. A repeated CNN
fit on one dataset, using identical observations and folds on two hardware platforms, produced
per-window score correlations of 0.9903 and fold-AUROC differences no larger than 0.0007. This
experiment combines hardware and initialisation effects and does not establish complete neural
reproducibility. All downstream analyses are deterministic given the released scores.

The positive sets in the two composition-matched arms were nearly, but not perfectly,
identical because a positive was discarded when its matcher found no acceptable negative. Their
median Jaccard similarity was 0.9972. Restricting both arms to the common positive set removed
0.23\% of positives and changed the contrast from $+0.0398$ to $+0.0401$, preserving its sign
in the same 88 of 94 datasets. Because matching failures are composition-dependent, this
intersection analysis is a sensitivity check rather than the primary estimand.

Finally, the released ENCODE peaks were already thresholded, and I applied no additional
significance filter. Fixed 101-nt windows discard peak width and may omit the crosslink site for
wide peaks. Negatives were not filtered by gene expression, and transcript region was assigned
from the window midpoint under a fixed annotation precedence. These choices may affect the
absolute measurements even though the same rules were applied across the paired protocols where
applicable.

The repository's verification suite provides regression and provenance checks, not independent
confirmation of the original analysis. Most assertions compare reported values with committed
evidence from the same run; two checks independently reconstruct 285 AUROCs from per-window
scores and the headline contrast from raw sequence. The manuscript-number trace can identify an
unsupported value but cannot determine whether a valid value has been attached to the correct
claim. These safeguards reduce transcription and drift errors, but scientific validity still
depends on the design, assumptions and limitations described above.
